\section{Concluding Remarks and Outlook}

We set out to ask whether the type IIB string theory admits a 12d origin, and what such an origin would reveal about its \SLZ{} duality.
We now summarise what the preceding sections establish, assessed against the objectives set out in the introduction, before turning to open directions.

\subsection{Aims and what was achieved}

\begin{enumerate}
    \item The 12d interpretation of the axio-dilaton sector (Section~\ref{sec:axio-dilaton-12d}) is achieved. 
    The D7 and D(-1) backgrounds arise respectively as a 12d Kaluza-Klein monopole and pp-wave reduced on a torus, the type IIB \SLZ{} acts as large gauge transformations in 12d, and the reduction reproduces the type IIB Killing spinor equations. 
    This is the sector that most robustly admits a 12d interpretation.
    \item The extension beyond the axio-dilaton sector (Section~\ref{sec:covariant-unification}) is partial. The RR and NSNS form fields can be assembled into 12d objects, with the 10d self-duality of the 5-form realised as a 12d Hodge duality. 
    This is, however, a covariant repackaging of the 10d theory on a two-torus: not all 12d degrees of freedom are excited, and the mismatch between the SO(8) and SO(10) little-group content shows that a full 12d uplift of the field content is unavailable.
    \item The role of the D3 brane (Section~\ref{sec:d3-self-duality}) is clarified. The worldvolume electromagnetic \SLZ{} duality of the D3 brane is consistent only when accompanied by the bulk \SLZ{} transformation, tying the worldvolume and bulk dualities together, following \cite{Tseytlin:1996it}.
    The potential 12d interpretations of this phenomenon are discussed.
    \item The 12d interpretation is tested against the higher-derivative corrections (Section~\ref{sec:effective-actions}), and the limits of its validity are made precise. 
    At four points the 8-derivative effective action can be written in 12d-covariant form, but at five points a U(1)-preserving axio-dilaton term obstructs it. Whether this obstruction is genuine or is lifted at higher points is left as an open question.
\end{enumerate}

Taken together, the 12d interpretation is robust for the axio-dilaton sector but, beyond it, remains effective rather than fundamental.
This is consistent with the long-standing obstruction to a 10+2d supergravity: what emerges is a 12d-covariant organisation of the 10d theory rather than a dynamical 12d spacetime.

\subsection{Outlook}

Two directions stand out for future work. The first concerns the effective action: whether a more general 12d ansatz, going beyond the $t_8 t_8 \widehat{R}^4$ contraction, can accommodate the type IIB couplings, and whether the five-point obstruction persists at higher points. It is possible that, beginning at six points, further obstructive terms appear that conspire with the five-point terms of \cite{Liu:2025uqu} to restore a 12d-covariant form.
Confirming this would require deriving the axio-dilaton effective action up to eight points.
The alternative is that the obstruction is genuine, and the type IIB effective action cannot be written 12d-covariantly.

The second concerns the role of \SLZ{} in the context of the AdS/CFT correspondence.
Recently, it was shown that an M2 brane wrapping a circle at the boundary of the $\mathrm{AdS}_4 \times S^7/\mathbb{Z}_k$ background reproduces, via a one-loop computation of its worldvolume effective action, the subleading $1/N$ corrections to Wilson-loop observables in the dual ABJM theory \cite{Giombi:2023vzu, Tseytlin:2024euk}.
This raises the question of whether an analogous construction exists for the $\mathrm{AdS}_5 \times S^5$ background, whose dual $\mathcal{N}=4$ SYM theory carries its own \SLZ{} duality.
Addressing it would require identifying the type IIB counterparts of the $\mathrm{AdS}_4 \times S^7/\mathbb{Z}_k$ geometry and the M2 brane, and in particular a 12d interpretation of the D3 brane that sources $\mathrm{AdS}_5 \times S^5$.
A sharper understanding of the 12d structure of type IIB string theory may therefore bear on whether such a correspondence exists.

Recently, the 10d non-supersymmetric type 0 string theories have received renewed attention.
It has been argued that the type 0A theory arises as M-theory compactified on a ``figure-8'' circle $S^1 \wedge S^1$ \cite{Baykara:2026gem}, while the type 0B theory with gauge group $SO(16)^4$ has been suggested to arise as M-theory on $S^1 \wedge S^1 \times S^1 / \mathbb{Z}_2$ \cite{Altavista:2026evd}.
An extension of the web of dualities into the type 0 theories has also been explored \cite{Baykara:2026gem, Baykara:2026vdc}.
How a 12d perspective may fit into this picture remains an interesting open question.